\documentclass[11pt]{article}
\usepackage[margin=1in]{geometry}
\usepackage{amsmath,amssymb,amsthm}
\usepackage[T1]{fontenc}
\usepackage{lmodern}
\usepackage{xcolor}
\usepackage{listings}
\usepackage{booktabs}
\usepackage{longtable}
\usepackage{array}
\usepackage[hidelinks]{hyperref}
\usepackage{url}

\newtheorem{theorem}{Theorem}
\newtheorem{lemma}[theorem]{Lemma}
\newtheorem{corollary}[theorem]{Corollary}
\theoremstyle{definition}
\newtheorem{definition}[theorem]{Definition}
\theoremstyle{remark}
\newtheorem{remark}[theorem]{Remark}

% Code-like text (not checked by a name gate).
\newcommand{\lean}[1]{\texttt{#1}}
% A declaration of one of the four developments. Every \decl{...} is meant to be checked by a
% name gate (grep for a declaration of that name in the four source trees).
\DeclareUrlCommand\decl{\urlstyle{tt}\def\UrlBreaks{\do\_\do\.}}
% A Mathlib declaration (pinned revision 0df444a3); not one of ours.
\DeclareUrlCommand\mdecl{\urlstyle{tt}\def\UrlBreaks{\do\_\do\.}}
\DeclareUnicodeCharacter{211D}{\ensuremath{\mathbb{R}}}
\DeclareUnicodeCharacter{2102}{\ensuremath{\mathbb{C}}}
\newcommand{\RR}{\mathbb{R}}
\newcommand{\CC}{\mathbb{C}}
\newcommand{\QQ}{\mathbb{Q}}
% Outputs that a separate job is producing. Must all be filled in before submission.
\newcommand{\PH}[1]{\textcolor{red}{\textbf{[PLACEHOLDER:} #1\textbf{]}}}

% ---------------------------------------------------------------------------------------------
% PINS. All four developments are PUBLISHED (github.com/jpbald93/<name>) at the commits below
% (verified via the GitHub API, 7 Oct 2026). For all four, the .lean sources at those commits are
% byte-identical (SHA-256 / cmp), file for file, to the copies this paper was written from
% (the copies also hold an unpublished scratch/Ax*.lean each; cited nowhere).
% ---------------------------------------------------------------------------------------------
\newcommand{\commitfeuerbach}{\texttt{2bf41901548aed8b0e43d2f1bd10bb1a1d54ebf8}}
\newcommand{\commitbrahmagupta}{\texttt{c5c3ce014dc450ca212df76fade8948aa45421e0}}
\newcommand{\commitmiquel}{\texttt{401de13c0c4e51f67f87f7252f147048084c4a01}}
\newcommand{\commitmalfatti}{\texttt{d39519346e71c9af8c349490a4dac5388de4d8f2}}
% Plain-text forms used inside URLs.
\expandafter\def\csname commitof:feuerbach\endcsname{2bf41901548aed8b0e43d2f1bd10bb1a1d54ebf8}
\expandafter\def\csname commitof:brahmagupta\endcsname{c5c3ce014dc450ca212df76fade8948aa45421e0}
\expandafter\def\csname commitof:miquel\endcsname{401de13c0c4e51f67f87f7252f147048084c4a01}
\expandafter\def\csname commitof:malfatti\endcsname{d39519346e71c9af8c349490a4dac5388de4d8f2}
\newcommand{\commitof}[1]{\csname commitof:#1\endcsname}
\newcommand{\repo}[1]{https://github.com/jpbald93/#1}
% \src{repo}{Dir}{File}{n}: line n of Dir/File.lean in the given repository. It is a link for
% the published repositories (those with a commitof: entry above) and plain text otherwise.
\newcommand{\src}[4]{\ifcsname commitof:#1\endcsname
  \href{\repo{#1}/blob/\commitof{#1}/#2/#3.lean\#L#4}{\texttt{#3}:#4}\else\texttt{#3}:#4\fi}

\lstdefinelanguage{lean}{
  morekeywords={theorem,lemma,def,abbrev,noncomputable,let,by,fun,namespace,end,import},
  sensitive=true, morecomment=[l]{--}, morecomment=[s]{/-}{-/}}
\lstset{language=lean, basicstyle=\ttfamily\small, keywordstyle=\bfseries,
  columns=fullflexible, keepspaces=true, breaklines=true, frame=single, framerule=0.3pt,
  literate={≠}{{$\neq$}}1 {→}{{$\to$}}1 {≤}{{$\le$}}1 {ℂ}{{$\mathbb{C}$}}1
           {ℝ}{{$\mathbb{R}$}}1 {∀}{{$\forall$}}1 {∃}{{$\exists$}}1 {∧}{{$\wedge$}}1
           {√}{{$\surd$}}1 {₁}{{$_1$}}1 {₂}{{$_2$}}1 {₃}{{$_3$}}1}

\title{The Malfatti circles in Lean~4, with independent\\
certificate proofs of Feuerbach, Miquel and Brahmagupta}
\author{Josh Bald\\ \small Independent researcher, Ontario, Canada\\ \small\texttt{jpbald93@gmail.com}\\
\small ORCID: \texttt{0009-0002-1317-6489}
\thanks{\copyright~2026 Josh Bald. Licensed under CC BY 4.0 (\url{https://creativecommons.org/licenses/by/4.0/}). The accompanying Lean and Python code is licensed under the Apache License 2.0.}}
\date{7 October 2026}

\emergencystretch=3em
\begin{document}
\sloppy
\maketitle

\begin{abstract}
We formalise four classical theorems about circles in Lean~4 with Mathlib. Only one of them,
the Malfatti circle construction, is to our knowledge formalised here for the first time; the other
three have earlier formal proofs and are offered as independent formalisations. \emph{Feuerbach's
theorem} (number~29 on Wiedijk's list of 100 theorems): the nine-point circle of a triangle is
internally tangent to the incircle and externally tangent to the three excircles, stated with
Mathlib's own \mdecl{Affine.Simplex.insphere}, \mdecl{Affine.Simplex.exsphere},
\mdecl{Affine.Simplex.ninePointCircle} and tangency predicates, together with Euler's formula
$OI^2=R(R-2r)$, its excentral analogue and Euler's inequality $2r\le R$. \emph{Miquel's theorem}:
for points on the three (extended) sides of a triangle, none at a vertex, the three Miquel circles
pass through a common point. \emph{Brahmagupta's formula}: a convex cyclic quadrilateral with
sides $a,b,c,d$ has (shoelace) area $\sqrt{(s-a)(s-b)(s-c)(s-d)}$, the formal hypotheses also
admitting some collapsed configurations for which the formula holds trivially, derived from Mathlib's Ptolemy theorem
through the diagonal form of Bretschneider's formula, which we prove for every quadrilateral.
\emph{The Malfatti circles}: every triangle in the Euclidean plane has three circles of positive
radius, each tangent to the two sides through one vertex and any two externally tangent, with
Malfatti's radius formula $\rho_A=\frac{r}{2(s-a)}\bigl(s-r-(IB+IC-IA)\bigr)$; this is the
construction, \emph{not} Malfatti's conjecture that these circles maximise the total area, which
Goldberg disproved in 1967~\cite{Goldberg}.
In each case the core is an explicit polynomial identity closed by \lean{ring} or
\lean{linear\_combination}. The four libraries use no \lean{sorry} and depend only on Lean's three
standard axioms; each passes an axiom gate, a tamper test and a further battery of statement-level
checks. All four libraries were also rebuilt on a second machine of a different architecture
(arm64/macOS).
\textbf{Feuerbach's theorem is not new to Lean:} Weiyi Wang has a sorry-free standalone Lean
repository (January 2026) and an open Mathlib pull request (\#44236, created 26 September 2026).
Our Feuerbach development is an independent second formalisation, and we claim no priority for
it. \textbf{Miquel's theorem is not new to formal mathematics either:} it was formalised in
Isabelle/HOL in the Archive of Formal Proofs entry of Freitas Ramos, Barros Hulak and de~Queiroz
(August 2026), and earlier in Coq's HighSchoolGeometry library. Our Miquel statement differs in
assuming only the non-degeneracy conditions that make the three circles exist.
\textbf{Brahmagupta's formula is not new to Lean:} a complete, sorry-free Lean~4 proof is in the
ShadowBench repository (\texttt{epfl-lara/icml-26-lean-challenges}, commit of 12 June 2026). What
our Brahmagupta development adds is limited to statements we did not find there: the
inequality $K^2\le(s-a)(s-b)(s-c)(s-d)$ and the equivalence of equality with Ptolemy's equality,
both for arbitrary four points. For the Malfatti circle construction we found no earlier
formalisation; our searches are not exhaustive, and three of the four searches missed existing
work. The new contributions are therefore the Malfatti formalisation, a second Lean proof of
Feuerbach's theorem, a Miquel statement with fewer hypotheses than the earlier ones, and the
verification methodology with its limits.
\end{abstract}

\section{Introduction}

Circles are where elementary geometry becomes awkward to formalise. A statement such as ``the
three circles pass through a common point'' presupposes that the circles exist and are unique;
``the area of a cyclic quadrilateral'' presupposes an order of the vertices and a notion of area;
``the nine-point circle touches the incircle'' must decide what happens for an equilateral
triangle, where the two circles coincide. A formal statement has to settle each of these
questions, and the settlement is part of the result.

This paper continues an earlier one on Desargues, Monge, Simson, Napoleon and Lami~\cite{PaperA}
and uses the same method: each theorem is reduced to an identity in coordinates,
\[
  (\text{conclusion})\cdot(\text{non-degenerate factor}) \;=\; \sum_i c_i\cdot(\text{hypothesis}_i),
\]
with the multipliers $c_i$ written in the source, closed by a decision procedure for commutative
rings. We call such an identity a \emph{certificate}. The method is classical~\cite{Chou}; what a
formalisation adds is a translation layer to standard Mathlib notions, and an exact account of
which non-degeneracy hypotheses are needed and where.

\paragraph{Contributions.}
\begin{itemize}
\item \emph{Feuerbach's theorem} for a triangle in an arbitrary real inner product space (the
plane hypothesis of the headline statement is present but unused), with tangency in Mathlib's
sense, including the degenerate equilateral case; Euler's formula $OI^2=R(R-2r)$, the excentral
formula $OI_i^2=R(R+2r_i)$, Euler's inequality $2r\le R$, and the distances
$|NI|=R/2-r$, $|NI_i|=R/2+r_i$ (Section~\ref{sec:feuerbach}). This is \emph{not} the first Lean
proof of Feuerbach's theorem (Section~\ref{sec:related}).
\item \emph{Miquel's theorem} in $\CC$, on the extended sides, under exactly the hypotheses that
make the three circles exist and be unique, in a ``for every circle through these three points''
form and in a form that constructs the three centres (Section~\ref{sec:miquel}). This is
\emph{not} the first formal proof of Miquel's theorem; it is the first in Lean that we know of, and
its hypotheses are weaker than those of the Isabelle version (Section~\ref{sec:related}).
\item \emph{Brahmagupta's formula} for convex cyclic quadrilaterals in
$\mathrm{EuclideanSpace}\ \RR\ (\mathrm{Fin}\ 2)$, stated with Mathlib's \mdecl{Cospherical}
and \mdecl{Sbtw}, together with the diagonal Bretschneider identity for every quadrilateral, the
inequality $K^2\le(s-a)(s-b)(s-c)(s-d)$ for every quadrilateral, and the equivalence of equality
with Ptolemy's equality (Section~\ref{sec:brahmagupta}). The area is the shoelace value
(Section~\ref{sec:limits}). This is \emph{not} the first Lean proof of Brahmagupta's formula: a
complete proof, by essentially the same route, is in the ShadowBench repository
(Section~\ref{sec:related}). The only parts we did not find there are the two all-four-point
statements, the inequality and the Ptolemy equivalence.
\item The \emph{Malfatti circles} for every triangle in $\mathrm{EuclideanSpace}\ \RR\
(\mathrm{Fin}\ 2)$: existence of three circles of positive radius, circle $i$ tangent (Mathlib's
\mdecl{IsTangent}) to the two side lines through vertex $i$ with its centre strictly on the
interior side of each, any two externally tangent (\mdecl{IsExtTangent}), and Malfatti's radius
formula (Section~\ref{sec:malfatti}). This is the classical construction, not Malfatti's
maximal-area conjecture, which is false.
\item The verification record of the four developments: gate, tamper tests, the house battery,
non-vacuity witnesses, definition bridges and, for all four, a cross-platform rebuild on
arm64/macOS (Section~\ref{sec:verify}), together with the limits of these checks.
\end{itemize}

\paragraph{Why a paper and not only Mathlib pull requests.} The code alone does not record three
things that this paper does. First, the hypothesis analysis: which non-degeneracy conditions each
statement needs and why, how they compare with earlier formalisations (Miquel against the Isabelle
version, Feuerbach's equilateral case), and which hypotheses turned out to be unused. Second, the
passage from a polynomial certificate to a usable geometric statement in Mathlib's vocabulary,
which is where most of the work and most of the design choices lie (the coincident-intersection
branch of Miquel, the quarter-angle parametrisation of Malfatti, shoelace area for Brahmagupta).
Third, the verification methodology together with an honest account of its limits, including what
the mutation and satisfiability layers did \emph{not} cover. The developments themselves remain
candidates for Mathlib; for Feuerbach that route runs through Wang's open pull request.
In short, of the four theorems only the Malfatti construction is new as a formalisation;
Feuerbach and Brahmagupta are second Lean proofs, and Miquel is a further formal proof
distinguished by its hypotheses.

\subsection{Availability}\label{sec:avail}
All four developments are published:
\begin{itemize}
\item \lean{Feuerbach}: \url{https://github.com/jpbald93/feuerbach}, commit \commitfeuerbach;
\item \lean{Miquel}: \url{https://github.com/jpbald93/miquel}, commit \commitmiquel;
\item \lean{Brahmagupta}: \url{https://github.com/jpbald93/brahmagupta}, commit \commitbrahmagupta;
\item \lean{Malfatti}: \url{https://github.com/jpbald93/malfatti}, commit \commitmalfatti.
\end{itemize}
The code in all four repositories (Lean sources, gate and tamper scripts, and the numerical
checkers) is released under the Apache License 2.0, the licence used by Mathlib; this paper is
released under CC BY 4.0.
The line references in Table~\ref{tab:thms} are links to these pinned commits. For
all four repositories every \path{.lean} file at these commits is byte-identical to its counterpart
in the copies from which this paper was written and on which the checks of
Section~\ref{sec:verify} were run. (Those working copies additionally hold one unpublished
scratch file each, \path{scratch/Ax*.lean}, used only to print axiom reports; it is cited nowhere.) The toolchain is Lean \texttt{v4.33.1}~\cite{Lean4} with Mathlib~\cite{Mathlib}
revision \texttt{0df444a3} (tag \texttt{v4.33.1}), pinned in each \path{lake-manifest.json}. The
libraries, counted on those copies, are \lean{Feuerbach} (five files and an import-only root file,
630 lines in total), \lean{Miquel} (six files and a root, 1{,}036 lines, of which 346 are the
definition bridges), \lean{Brahmagupta} (four files and a root, 363 lines) and \lean{Malfatti}
(five files and a root, 984 lines, of which 202 are the definition bridges and 277 the
non-vacuity witnesses). All declarations are in the namespace of the library name.

{\small
\setlength{\LTcapwidth}{\textwidth}
\begin{longtable}{@{}>{\raggedright\arraybackslash}p{0.27\textwidth}
  >{\raggedright\arraybackslash}p{0.33\textwidth}
  >{\raggedright\arraybackslash}p{0.12\textwidth}
  >{\raggedright\arraybackslash}p{0.20\textwidth}@{}}
\caption[Informal statements and their formal counterparts]{Informal statements and their formal counterparts. Line numbers refer to the first line
of the declaration. For Feuerbach, $t$ is an \texttt{Affine.Triangle} in a real inner product
space (no dimension hypothesis unless stated).}\label{tab:thms}\\
\toprule
Informal statement & Lean name & Source & Hypotheses \\ \midrule
\endfirsthead
\toprule
Informal statement & Lean name & Source & Hypotheses \\ \midrule
\endhead
\bottomrule
\endfoot
\multicolumn{4}{@{}l}{\emph{Feuerbach (Section~\ref{sec:feuerbach})}}\\
Circumradius relation $4\,\mathrm{heron}\cdot R^2=a^2b^2c^2$ & \decl{Feuerbach.circ_relation} & \src{feuerbach}{Feuerbach}{Algebra}{22} & three equidistance equations \\
Certificate $|N-I|^2=(R/2-kT)^2$ & \decl{Feuerbach.nine_in} & \src{feuerbach}{Feuerbach}{Algebra}{28} & frame equations, $k^2=\mathrm{heron}$, $2kR=abc$, $T(a+b+c)=1$ \\
$h_i^2\cdot(\text{opposite side})^2=\mathrm{heron}$ & \decl{Feuerbach.height_sq_mul_sq} & \src{feuerbach}{Feuerbach}{Geometry}{43} & $t$ \\
Euler: $OI^2=R(R-2r)$ & \decl{Feuerbach.dist_circumcenter_incenter_sq} & \src{feuerbach}{Feuerbach}{Main}{100} & $t$ \\
Excentral Euler: $OI_i^2=R(R+2r_i)$ & \decl{Feuerbach.dist_circumcenter_excenter_sq} & \src{feuerbach}{Feuerbach}{Main}{105} & $t$, $i$ \\
Euler's inequality $2r\le R$ & \decl{Feuerbach.two_mul_inradius_le_circumradius} & \src{feuerbach}{Feuerbach}{Main}{111} & $t$ \\
$|NI|=R/2-r$ & \decl{Feuerbach.dist_ninePointCircle_center_incenter} & \src{feuerbach}{Feuerbach}{Main}{119} & $t$ \\
$|NI_i|=R/2+r_i$ & \decl{Feuerbach.dist_ninePointCircle_center_excenter} & \src{feuerbach}{Feuerbach}{Main}{126} & $t$, $i$ \\
Incircle internally tangent & \decl{Feuerbach.isIntTangent_insphere_ninePointCircle} & \src{feuerbach}{Feuerbach}{Main}{143} & $t$ \\
Excircles externally tangent & \decl{Feuerbach.isExtTangent_exsphere_ninePointCircle} & \src{feuerbach}{Feuerbach}{Main}{154} & $t$, $i$ \\
Feuerbach's theorem (Thm.~\ref{thm:feuerbach}) & \decl{Feuerbach.feuerbach} & \src{feuerbach}{Feuerbach}{Main}{163} & $t$; \lean{Fact (finrank ℝ V = 2)} (unused) \\
Non-vacuity on the 3-4-5 triangle & \decl{Feuerbach.feuerbach_nonvacuous} & \src{feuerbach}{Feuerbach}{Main}{171} & none \\
Bridge: $\mathrm{qf}=\|pu+qv\|^2$ & \decl{Feuerbach.qf_eq_norm_sq} & \src{feuerbach}{Feuerbach}{Bridge}{26} & none \\
Bridge: $\mathrm{heron}=(bc\sin A)^2$ & \decl{Feuerbach.heron_eq_sq_two_mul_area} & \src{feuerbach}{Feuerbach}{Bridge}{36} & none \\
\multicolumn{4}{@{}l}{\emph{Miquel (Section~\ref{sec:miquel}); points in $\CC$}}\\
Orientation $\mathrm{Dp}$ & \decl{Miquel.Dp} & \src{miquel}{Miquel}{Frame}{26} & definition \\
Concyclicity determinant $Q$ & \decl{Miquel.Q} & \src{miquel}{Miquel}{Basic}{23} & definition \\
Certificates ($x$- and $y$-part) & \decl{Miquel.cert_x}, \decl{Miquel.cert_y} & \src{miquel}{Miquel}{Coeffs}{91}, \src{miquel}{Miquel}{Coeffs}{98} & none \\
Tangency and separation identities & \decl{Miquel.tang}, \decl{Miquel.sep} & \src{miquel}{Miquel}{Coeffs}{106}, \src{miquel}{Miquel}{Coeffs}{113} & none \\
Miquel in the frame & \decl{Miquel.frame} & \src{miquel}{Miquel}{Frame}{167} & $\Delta\ne0$, $n\ne0,1$, $t\ne0$, $u\ne0$ \\
Miquel's theorem (Thm.~\ref{thm:miquel}) & \decl{Miquel.miquel} & \src{miquel}{Miquel}{Main}{54} & $\mathrm{Dp}(A,B,C)\ne0$; $n,t,u\notin\{0,1\}$ \\
Miquel, concyclic form & \decl{Miquel.miquel_concyclic} & \src{miquel}{Miquel}{Main}{97} & same \\
Example (midpoints) & \decl{Miquel.miquel_example} & \src{miquel}{Miquel}{Main}{136} & none \\
Negative control & \decl{Miquel.miquel_negative_control} & \src{miquel}{Miquel}{Main}{155} & none \\
All ten hypotheses met & \decl{Miquel.miquel_hyps_example} & \src{miquel}{Miquel}{Witness}{27} & none \\
Witness for \lean{core}: second intersection is $(1,1)$ & \decl{Miquel.core_example_point} & \src{miquel}{Miquel}{Witness}{111} & none \\
Bridge: $\mathrm{Dp}=0\iff$ \mdecl{Collinear} & \decl{Miquel.Dp_eq_zero_iff_collinear} & \src{miquel}{Miquel}{Bridge}{44} & none \\
Bridge: $Q=\mathrm{Dp}\cdot$ \mdecl{Sphere.power} & \decl{Miquel.Q_eq_Dp_mul_power} & \src{miquel}{Miquel}{Bridge}{92} & $|R-O|=|S-O|=|P-O|$ \\
\multicolumn{4}{@{}l}{\emph{Brahmagupta (Section~\ref{sec:brahmagupta}); points in $\mathrm{EuclideanSpace}\ \RR\ (\mathrm{Fin}\ 2)$}}\\
Signed (shoelace) area & \decl{Brahmagupta.signedArea} & \src{brahmagupta}{Brahmagupta}{Basic}{25} & definition \\
Area $K=|\text{signed area}|$ & \decl{Brahmagupta.area} & \src{brahmagupta}{Brahmagupta}{Basic}{29} & definition \\
Equals the shoelace sum & \decl{Brahmagupta.signedArea_eq_shoelace} & \src{brahmagupta}{Brahmagupta}{Basic}{31} & none \\
Diagonal Bretschneider $16K^2=4p^2q^2-(b^2+d^2-a^2-c^2)^2$ & \decl{Brahmagupta.sixteen_area_sq} & \src{brahmagupta}{Brahmagupta}{Basic}{47} & none (any 4 points) \\
Algebraic identity & \decl{Brahmagupta.brahmagupta_identity} & \src{brahmagupta}{Brahmagupta}{Basic}{55} & none \\
$K^2\le\prod(s-\cdot)$ & \decl{Brahmagupta.area_sq_le} & \src{brahmagupta}{Brahmagupta}{Basic}{65} & none (any 4 points) \\
$K^2=\prod(s-\cdot)\iff pq=ac+bd$ & \decl{Brahmagupta.area_sq_eq_iff_ptolemy} & \src{brahmagupta}{Brahmagupta}{Cyclic}{21} & none (any 4 points) \\
Brahmagupta, squared & \decl{Brahmagupta.brahmagupta_sq} & \src{brahmagupta}{Brahmagupta}{Cyclic}{47} & \mdecl{Cospherical}; diagonals cross (\mdecl{Sbtw}) \\
Brahmagupta's formula (Thm.~\ref{thm:brahmagupta}) & \decl{Brahmagupta.brahmagupta} & \src{brahmagupta}{Brahmagupta}{Cyclic}{63} & same \\
Contrapositive test for cyclicity & \decl{Brahmagupta.not_cospherical_of_area_sq_ne} & \src{brahmagupta}{Brahmagupta}{Cyclic}{76} & diagonals cross; $K^2\ne\prod(s-\cdot)$ \\
Unit square: all hypotheses & \decl{Brahmagupta.unitSquare_example} & \src{brahmagupta}{Brahmagupta}{Examples}{66} & none \\
Rhombus is not cyclic & \decl{Brahmagupta.rhombus_not_cospherical} & \src{brahmagupta}{Brahmagupta}{Examples}{100} & none \\
Bridge: $K=\tfrac12\,pq\sin\theta$ & \decl{Brahmagupta.area_eq_half_mul_sin_angle} & \src{brahmagupta}{Brahmagupta}{Bridge}{28} & none \\
\multicolumn{4}{@{}l}{\emph{Malfatti (Section~\ref{sec:malfatti}); $t$ a triangle; headlines in $\mathrm{EuclideanSpace}\ \RR\ (\mathrm{Fin}\ 2)$}}\\
Tangent relation $K(p,q,w)=0$ from Heron & \decl{Malfatti.K_eq_zero_of_heron} & \src{malfatti}{Malfatti}{Algebra}{31} & $r>0$; $p,q,w\in(0,1)$; Heron \\
Radius formula, parametrised & \decl{Malfatti.radius_eq} & \src{malfatti}{Malfatti}{Algebra}{52} & $p\in(0,1)$, $q,w>0$, $K=0$ ($r>0$ unused) \\
Pairwise tangency identity & \decl{Malfatti.dist_identity} & \src{malfatti}{Malfatti}{Algebra}{69} & $p,q>0$, $K=0$ ($p<1$ unused) \\
Heron $r^2s=(s-a)(s-b)(s-c)$ & \decl{Malfatti.heron} & \src{malfatti}{Malfatti}{Geometry}{134} & $t$; \lean{finrank ℝ V = 2} \\
The circles (definition) & \decl{Malfatti.malfattiCircle} & \src{malfatti}{Malfatti}{Geometry}{34} & definition \\
Side tangency, interior side & \decl{Malfatti.malfattiCircle_tangent_side} & \src{malfatti}{Malfatti}{Geometry}{240} & $t$; \lean{finrank ℝ V = 2}; $i,j,k$ distinct \\
Pairwise external tangency & \decl{Malfatti.malfattiCircle_isExtTangent} & \src{malfatti}{Malfatti}{Geometry}{280} & $t$; \lean{finrank ℝ V = 2}; $i\ne j$ \\
Radius formula & \decl{Malfatti.malfattiCircle_radius_formula} & \src{malfatti}{Malfatti}{Geometry}{298} & $t$; \lean{finrank ℝ V = 2} \\
Existence (Thm.~\ref{thm:malfatti}) & \decl{Malfatti.malfatti_circles_exist} & \src{malfatti}{Malfatti}{Main}{15} & $t$ \\
Existence with radius formula & \decl{Malfatti.malfatti_circles_formula} & \src{malfatti}{Malfatti}{Main}{31} & $t$ \\
Instance at $(0,0),(1,0),(0,1)$ & \decl{Malfatti.malfatti_circles_formula_triW} & \src{malfatti}{Malfatti}{Witness}{97} & none \\
Bridge: $K(\tan a,\tan b,\tan c)=0\iff a+b+c=\pi/4$ & \decl{Malfatti.K_bridge} & \src{malfatti}{Malfatti}{Bridge}{50} & $a,b,c\in(0,\pi/4)$ \\
Bridge: $D(r,\tan\theta)=r/\sin2\theta$ & \decl{Malfatti.D_bridge} & \src{malfatti}{Malfatti}{Bridge}{38} & $\theta\in(0,\pi/2)$ \\
Bridge: $\mathrm{tparam}=\tan(A_i/4)$ & \decl{Malfatti.tparam_bridge} & \src{malfatti}{Malfatti}{Bridge}{115} & $t$; \lean{finrank ℝ V = 2} \\
Bridge: $r\mu(\ldots)$ is the classical radius & \decl{Malfatti.mu_bridge} & \src{malfatti}{Malfatti}{Bridge}{162} & $t$; \lean{finrank ℝ V = 2} \\
\end{longtable}
}

\section{The method}\label{sec:method}

The four proofs follow the three steps of~\cite{PaperA}.
\begin{enumerate}
\item \emph{Coordinates.} Miquel's theorem is stated for points of $\CC$ and Brahmagupta's for
points of $\mathrm{EuclideanSpace}\ \RR\ (\mathrm{Fin}\ 2)$; in both, the real and imaginary parts
(or the two coordinates) are the variables. Feuerbach's theorem is stated for an
\mdecl{Affine.Triangle} in an arbitrary real inner product space, and no isometry to $\RR^2$ or
$\CC$ is used. Instead the argument works in the affine frame $(p_0;\,p_1-p_0,\,p_2-p_0)$ of the
triangle, with the three side lengths $a,b,c$ as real parameters: every Gram entry is a polynomial
in $a,b,c$, and square roots never appear because the side lengths are themselves variables.
The Malfatti development likewise works with a triangle in a general 2-dimensional real inner
product space; its variables are the inradius $r$ and the tangents $p,q,w$ of the quarter angles,
so that every length is a rational function of them.
\item \emph{Certificate.} After clearing denominators the conclusion is an explicit combination of
hypothesis polynomials. The multipliers are computed outside Lean (with SymPy) and written into the
source; Lean checks them with \lean{ring} or \lean{linear\_combination}. No proof uses
\lean{native\_decide} or \lean{polyrith}; \lean{decide} appears only for finite facts about
$\mathrm{Fin}\ 3$.
\item \emph{Translation.} The coordinate statement is connected with Mathlib's notions:
\mdecl{EuclideanGeometry.Sphere.IsIntTangent}, \mdecl{EuclideanGeometry.Sphere.IsExtTangent},
\mdecl{Affine.Simplex.circumcenter}, \mdecl{Affine.Simplex.incenter} and
\mdecl{Affine.Simplex.excenter} for Feuerbach; \mdecl{dist} for Miquel, with bridges to
\mdecl{Collinear} and \mdecl{EuclideanGeometry.Sphere.power}; \mdecl{EuclideanGeometry.Cospherical},
\mdecl{Sbtw} and Mathlib's Ptolemy theorem for Brahmagupta; \mdecl{Affine.Simplex.incenter},
\mdecl{Affine.Simplex.inradius}, \mdecl{Affine.Simplex.touchpoint},
\mdecl{EuclideanGeometry.Sphere.IsTangent}, \mdecl{EuclideanGeometry.Sphere.IsExtTangent} and
\mdecl{AffineSubspace.SSameSide} for Malfatti.
\end{enumerate}
The certificates differ greatly in size. Brahmagupta's are closed by \lean{ring} in a few lines.
Miquel's are linear in the unknown point but of moderate size. Feuerbach's \lean{nine\_*} and
\lean{euler\_*} certificates are long (a few hundred monomials each) but each compiles in under a
second; a clean build of the whole Feuerbach library took 14~s on the build machine.

\section{Feuerbach's theorem}\label{sec:feuerbach}

\paragraph{Statement.} Let $V$ be a real inner product space, $P$ a metric affine space over $V$,
and $t$ an \mdecl{Affine.Triangle}~$\RR$~$P$ (three affinely independent points). Mathlib defines
the circumcentre $O$, circumradius $R$, incentre $I$ and inradius $r$, the excentres $I_i$ and
exradii $r_i$ (as \lean{t.excenter \{i\}}, \lean{t.exradius \{i\}}), the insphere, exspheres and
the nine-point circle with centre $N$, and the predicates \mdecl{IsIntTangent} and
\mdecl{IsExtTangent} for spheres.

\begin{theorem}[\texttt{Feuerbach.feuerbach}, \src{feuerbach}{Feuerbach}{Main}{163}]\label{thm:feuerbach}
For every triangle $t$, the incircle of $t$ is internally tangent to the nine-point circle, and for
each $i\in\mathrm{Fin}\ 3$ the excircle opposite vertex $i$ is externally tangent to the nine-point
circle:
\begin{lstlisting}
theorem feuerbach [Fact (Module.finrank ℝ V = 2)] (t : Triangle ℝ P) :
    t.insphere.IsIntTangent t.ninePointCircle ∧
      ∀ i : Fin 3, (t.exsphere {i}).IsExtTangent t.ninePointCircle
\end{lstlisting}
\end{theorem}

The two conjuncts are proved separately as \decl{Feuerbach.isIntTangent_insphere_ninePointCircle}
and \decl{Feuerbach.isExtTangent_exsphere_ninePointCircle}, with \emph{no} dimension hypothesis:
everything the argument uses lies in the plane of the triangle. The hypothesis
\lean{Fact (finrank ℝ V = 2)} appears in \decl{Feuerbach.feuerbach} only because it is part of the
target statement, and it is unused. Non-degeneracy is free, since an \mdecl{Affine.Triangle} is
affinely independent by definition. The only extra ingredient for the internal case is
\lean{Nontrivial V}, which \decl{Feuerbach.nontrivial_of_triangle} derives from the triangle.

\paragraph{The equilateral case.} For an equilateral triangle $R/2-r=0$, so the incircle and the
nine-point circle coincide. Mathlib's \mdecl{IsIntTangent} includes this degenerate same-sphere
case, and the proof goes through \mdecl{EuclideanGeometry.Sphere.isIntTangent_iff_dist_center},
which handles it; no case split is needed. The numerical script confirms $|NI|=0=R/2-r$ for an
equilateral triangle.

\paragraph{Euler's formula and inequality.} On the way we prove
\[
 OI^2=R(R-2r),\qquad OI_i^2=R(R+2r_i),\qquad 2r\le R
\]
(\decl{Feuerbach.dist_circumcenter_incenter_sq}, \decl{Feuerbach.dist_circumcenter_excenter_sq},
\decl{Feuerbach.two_mul_inradius_le_circumradius}, \src{feuerbach}{Feuerbach}{Main}{100}--111).
A search of the pinned Mathlib for \lean{inradius} together with \lean{circumradius} found neither
Euler's formula nor the inequality, so they are new relative to that revision; Wang's repository
and pull request (Section~\ref{sec:related}) also prove versions of them.

\paragraph{Proof architecture.} The development has three layers.
\begin{itemize}
\item \path{Feuerbach/Algebra.lean} is pure real algebra. With $u=p_1-p_0$, $v=p_2-p_0$ and side
lengths $a=|p_1p_2|$, $b=|p_0p_2|$, $c=|p_0p_1|$, the quadratic form
$\mathrm{qf}(a,b,c,p,q)=p^2c^2+pq(b^2+c^2-a^2)+q^2b^2$ is $\|pu+qv\|^2$, and
$\mathrm{heron}(a,b,c)=(2a^2b^2+2b^2c^2+2c^2a^2-a^4-b^4-c^4)/4$ is the square of twice the area.
The circumcentre is $p_0+xu+yv$; $k$ is twice the area; $T=1/(\pm a\pm b\pm c)$, so that $kT$ is
the in- or exradius. The lemmas \decl{Feuerbach.circ_relation} ($4\,\mathrm{heron}\cdot
R^2=a^2b^2c^2$), \decl{Feuerbach.nine_in}, \decl{Feuerbach.nine_ex0}--\lean{2} (giving
$|N-I_\varepsilon|^2=(R/2\mp kT)^2$) and \decl{Feuerbach.euler_in}, \decl{Feuerbach.euler_ex0}--\lean{2}
(giving $|O-I_\varepsilon|^2=R^2\mp2R\,kT$) are each one \lean{linear\_combination}. The
multipliers were generated by solving the $T$-free homogenised identity by linear algebra over an
ansatz and adding the multiple of $T\cdot D-1$ by hand.
\item \path{Feuerbach/Geometry.lean} connects Mathlib's definitions with this algebra. The key
step is \decl{Feuerbach.height_sq_mul_sq}: $h_i^2\cdot(\text{opposite side})^2=\mathrm{heron}$, a
Gram-determinant identity (\decl{Feuerbach.gram_of_foot}) applied at the altitude foot. It turns
Mathlib's excentre weights, which are $\pm1/h_i$, into weights proportional to $\pm$ the side
lengths. \decl{Feuerbach.frame_data} expresses $O$ and $N$ in the frame and derives $2kR=abc$;
\decl{Feuerbach.excenter_data} gives the excentre coordinates and $r_\varepsilon=k/(\pm a\pm b\pm c)$,
taking positivity of the weight sum from Mathlib.
\item \path{Feuerbach/Main.lean} instantiates a generic lemma \decl{Feuerbach.core} at the incentre
and the three excentres. Euler's formula gives $2r\le R$, which fixes the sign of the square root
and yields $|NI|=R/2-r$; then Mathlib's \mdecl{EuclideanGeometry.Sphere.isIntTangent_iff_dist_center}
and \mdecl{EuclideanGeometry.Sphere.isExtTangent_iff_dist_center} give tangency.
\end{itemize}

\paragraph{Scope.} Only the three ``true'' excircles, with sign set $\{i\}$, are covered. For a
triangle the sign sets of size two and three give the same spheres as those of size one and
zero; Wang's version is stated for every sign set other than $\emptyset$ and the full set, a mild
generalisation that we have not added.

\paragraph{Numerical check.} The script \path{code/check_feuerbach.py} (not part of the proof)
computes $O$, $H$, $N=(O+H)/2$, $I$, $I_a$, $I_b$, $I_c$ from first principles and checks the four
distance identities: on 19{,}998 random triangles the worst relative error is $9.97\times10^{-8}$
(for badly conditioned triangles); isosceles and equilateral cases give errors of order $10^{-15}$;
exact SymPy checks on the 3-4-5 triangle, the 13-14-15 triangle and the triangle $(0,0),(3,0),(1,2)$
give residual~$0$ for the incircle and all three excircles.

\section{Miquel's theorem}\label{sec:miquel}

\paragraph{Model.} The plane is $\CC$ with Mathlib's \mdecl{dist}. For points $a,b,c$ let
$\mathrm{Dp}(a,b,c)$ (\decl{Miquel.Dp}) be minus the cross product of $b-a$ and $c-a$, so that
$\mathrm{Dp}(a,b,c)\ne0$ exactly when $a,b,c$ are not collinear
(\decl{Miquel.Dp_eq_zero_iff_collinear}). Points on the sides are given by real parameters, and the
\emph{extended} sides are allowed:
\[
F=A+n(B-A),\qquad E=A+t(C-A),\qquad D=B+u(C-B).
\]
The hypotheses $n,t,u\notin\{0,1\}$ say that $D,E,F$ are not vertices; together with
$\mathrm{Dp}(A,B,C)\ne0$ they are exactly what makes the triples $(A,E,F)$, $(B,F,D)$, $(C,D,E)$
non-collinear, so that each of the three circles exists and is unique.

\begin{theorem}[\texttt{Miquel.miquel}, \src{miquel}{Miquel}{Main}{54}]\label{thm:miquel}
Under these hypotheses there is $M\in\CC$ such that, for every $O\in\CC$ and $r\in\RR$: if
$A,E,F$ are at distance $r$ from $O$ then so is $M$; if $B,F,D$ are at distance $r$ from $O$ then
so is $M$; and if $C,D,E$ are at distance $r$ from $O$ then so is $M$.
\end{theorem}

The ``for every circle'' form is meaningful only because the three circles exist; the companion
\decl{Miquel.miquel_concyclic} (\src{miquel}{Miquel}{Main}{97}) states, under the same hypotheses,
that there are $O_1,O_2,O_3,M$ with $E,F,M$ on the circle through $A$ centred at $O_1$, $F,D,M$ on
the circle through $B$ centred at $O_2$, and $D,E,M$ on the circle through $C$ centred at $O_3$.
Miquel's theorem has earlier formal proofs, in Isabelle/HOL and in Coq; what distinguishes this
statement is that it assumes only the non-degeneracy conditions above (Section~\ref{sec:related}).

\paragraph{Certificate.} In a translation-invariant frame with $B=A+w$, $C=A+c$ and $M=F+(x,y)$,
the three circle conditions are polynomials $q_1,q_2,q_3$ in $(x,y)$ (\decl{Miquel.q1}, \decl{Miquel.q2},
\decl{Miquel.q3}). A plain identity $q_3=P\,q_1+Q\,q_2$ with linear $P,Q$ is impossible: $q_1$ and
$q_2$ have no constant term, since both circles pass through $F$, while $q_3$ does. Multiplying by
$x$ or $y$ repairs this. With $K=n\,t\,u\,(n-1)\,|w|^2$ (\decl{Miquel.KK}) the certificates are
\[
K\,x\,q_3=P_1q_1+P_2q_2,\qquad K\,y\,q_3=P_1'q_1+P_2'q_2
\]
(\decl{Miquel.cert_x}, \decl{Miquel.cert_y}), with $P_i,P_i'$ linear in $(x,y)$ and computed by
SymPy. Two further identities complete the argument: \decl{Miquel.tang} handles the case $M=F$
(the constant term of $q_3$ is a multiple of the cross product of the linear parts of $q_1,q_2$),
and \decl{Miquel.sep} shows that the radical-axis normal is non-zero. All are closed by
\lean{ring}. The second intersection point of the first two circles is given explicitly
(\decl{Miquel.rpx}, \decl{Miquel.rpy}); \decl{Miquel.rp_on} shows that it lies on both, and the
abstract lemma \decl{Miquel.core} uses the certificates to put it on the third, giving
\decl{Miquel.frame}. The lemma \decl{Miquel.bridge} identifies the concyclicity determinant
\decl{Miquel.Q} and $\mathrm{Dp}$ with the frame polynomials, again by \lean{ring}.

\paragraph{Examples.} \decl{Miquel.miquel_example} applies the theorem to $A=0$, $B=4$,
$C=1+3i$ with $D,E,F$ the midpoints, discharging all hypotheses. \decl{Miquel.miquel_negative_control}
shows that in the same configuration $F$, which lies on the first two circles, lies on no circle
through $C,D,E$, so membership of the third circle is a real constraint; numerically $F$ is at
distance about $0.94$ from that circle.

\paragraph{Numerical check.} \path{code/check_miquel_formal.py} checks the exact formal
configuration, extended sides included, on 20{,}000 random cases: the worst relative residual is
$2.3\times10^{-11}$, while the negative control has residual at least $2.7\times10^{-5}$.

\section{Brahmagupta's formula}\label{sec:brahmagupta}

\paragraph{Area.} Points are elements of \lean{Pt := EuclideanSpace ℝ (Fin 2)}. For a
quadrilateral $ABCD$ we define
\begin{align*}
\mathrm{signedArea}(A,B,C,D)&=\tfrac12\big((C_0-A_0)(D_1-B_1)-(C_1-A_1)(D_0-B_0)\big),\\
K=\mathrm{area}(A,B,C,D)&=|\mathrm{signedArea}(A,B,C,D)|,
\end{align*}
half the cross product of the diagonals; \decl{Brahmagupta.signedArea_eq_shoelace} proves that
this equals the shoelace sum $\tfrac12\sum(x_iy_{i+1}-x_{i+1}y_i)$. For a simple (non-self-intersecting)
quadrilateral this is the enclosed area. For a crossed quadrilateral it is the difference of the
two lobes, which is why the main theorems assume convexity (Section~\ref{sec:limits}).

\paragraph{Three steps.} Write $a=AB$, $b=BC$, $c=CD$, $d=DA$, $p=AC$, $q=BD$ and $s=(a+b+c+d)/2$.
\begin{enumerate}
\item For \emph{every} quadrilateral, $16K^2=4p^2q^2-(b^2+d^2-a^2-c^2)^2$
(\decl{Brahmagupta.sixteen_area_sq}): the diagonal form of Bretschneider's formula, proved by
\lean{ring} after expanding distances into coordinates.
\item As a polynomial identity,
$16(s-a)(s-b)(s-c)(s-d)=4(ac+bd)^2-(b^2+d^2-a^2-c^2)^2$ (\decl{Brahmagupta.brahmagupta_identity}).
\item Hence, for every quadrilateral, $K^2=(s-a)(s-b)(s-c)(s-d)$ holds if and only if Ptolemy's
equality $pq=ac+bd$ holds (\decl{Brahmagupta.area_sq_eq_iff_ptolemy}), and Ptolemy's inequality
in Mathlib (\mdecl{EuclideanGeometry.mul_dist_le_mul_dist_add_mul_dist}) gives
$K^2\le(s-a)(s-b)(s-c)(s-d)$ for every quadrilateral (\decl{Brahmagupta.area_sq_le}).
\end{enumerate}
Mathlib's Ptolemy theorem \mdecl{EuclideanGeometry.mul_dist_add_mul_dist_eq_mul_dist_of_cospherical}
supplies the equality for a convex cyclic quadrilateral.

\begin{theorem}[\texttt{Brahmagupta.brahmagupta}, \src{brahmagupta}{Brahmagupta}{Cyclic}{63}]\label{thm:brahmagupta}
Let $A,B,C,D,P$ be points of the Euclidean plane with $\{A,B,C,D\}$ \mdecl{Cospherical}, $P$
strictly between $A$ and $C$, and $P$ strictly between $B$ and $D$. Then
$K=\sqrt{(s-a)(s-b)(s-c)(s-d)}$.
\end{theorem}

\decl{Brahmagupta.brahmagupta_sq} is the squared form. The hypothesis
$\mathrm{Sbtw}\,A\,P\,C\wedge\mathrm{Sbtw}\,B\,P\,D$ says that the diagonals cross at a point
interior to both; it is what makes Mathlib's Ptolemy theorem apply. \mdecl{Sbtw} forces $A\ne C$
and $B\ne D$, but \emph{not} that the four vertices are distinct, so the hypotheses do not by
themselves describe a non-degenerate convex quadrilateral. For example, $A=B=(-1,0)$,
$C=D=(1,0)$, $P=(0,0)$ satisfies all three hypotheses (the points lie on the unit circle and $P$ is
strictly between $A$ and $C$ and between $B$ and $D$), yet the ``quadrilateral'' has collapsed to
a segment. The theorem remains true there, and trivially so: the shoelace area is $0$, the sides
are $0,2,0,2$, $s=2$ and $(s-a)(s-b)(s-c)(s-d)=0$. Such configurations are therefore admitted
rather than excluded; they are not counterexamples. When, in addition, the four vertices are
pairwise distinct, the two strict-between conditions do make $ABCD$ a convex quadrilateral with
the vertices in this order, all four sides are positive, and informally the shoelace value is
then its enclosed area; this last identification is not proved in Lean (Section~\ref{sec:limits}).
We have not added a corollary with an explicit distinctness hypothesis, since it would assume
more and conclude nothing new.

\paragraph{Relation to the earlier Lean proof.} Theorem~\ref{thm:brahmagupta} is not new. The
ShadowBench repository has a complete Lean~4 proof, \lean{brahmagupta\_formula}, of the same
statement (Section~\ref{sec:related}, \cite{ShadowBenchProof}), by essentially the same route:
diagonals crossing, Mathlib's Ptolemy theorem, a diagonal identity for the squared area proved
by \lean{ring} for arbitrary points, and the semiperimeter algebra. Its hypotheses are expressed
differently (an explicit centre and radius instead of \mdecl{Cospherical}; four orientation
signs instead of crossing diagonals; the area as the sum of two triangle areas, and $K$ supplied
through a hypothesis $K=\mathrm{quadrilateralArea}\ldots$), but these are presentational
differences, not mathematical ones. What our development adds beyond that proof is narrow:
\decl{Brahmagupta.area_sq_le} and \decl{Brahmagupta.area_sq_eq_iff_ptolemy}, which hold for
\emph{any} four points with no cyclicity or convexity hypothesis, and the contrapositive test
\decl{Brahmagupta.not_cospherical_of_area_sq_ne} derived from them, for which we found no earlier
formal proof; the statement in Mathlib's \mdecl{Cospherical}/\mdecl{Sbtw} vocabulary; the
definition bridge \decl{Brahmagupta.area_eq_half_mul_sin_angle}; the witnesses; and the
verification record of Section~\ref{sec:verify}. We did not compile the ShadowBench proof
locally (it targets Lean \texttt{v4.29.0}; an attempt under our \texttt{v4.33.1} toolchain timed
out), so our description of it rests on reading the source, which contains no \lean{sorry}.

\paragraph{Examples.} \decl{Brahmagupta.unitSquare_example} shows that the square with vertices
$(1,0),(0,1),(-1,0),(0,-1)$ satisfies all hypotheses and has area $2=(\sqrt2)^2$;
\decl{Brahmagupta.unitSquare_brahmagupta} applies both headline theorems to it.
\decl{Brahmagupta.rhombus_fails} shows that the convex rhombus $(1,0),(0,2),(-1,0),(0,-2)$ has
$K^2=16$ while $(s-a)(s-b)(s-c)(s-d)=25$, so by \decl{Brahmagupta.not_cospherical_of_area_sq_ne}
it is not cyclic (\decl{Brahmagupta.rhombus_not_cospherical}); the cyclic hypothesis cannot be
dropped. \path{code/check_brahmagupta.py} checks the formula on 10{,}000 random cyclic
quadrilaterals (relative error below $10^{-13}$) and the general identity on 10{,}000 random
quadrilaterals.

\section{The Malfatti circles}\label{sec:malfatti}

\paragraph{Construction, not conjecture.} Two different statements carry Malfatti's name, and
only the first is formalised here. (i) The \emph{Malfatti circles} of a triangle are three
circles, each tangent to the two sides through one vertex, any two of them externally tangent;
Malfatti (1803) gave a construction and the radius formula below. (ii) \emph{Malfatti's
conjecture} (or problem) asserted that these three circles have the largest total area of any
three non-overlapping circles in the triangle. The conjecture is \emph{false}: Lob and Richmond
observed that it already fails for the equilateral triangle~\cite{LobRichmond}, and
Goldberg~\cite{Goldberg} (1967) showed that the Malfatti circles are never optimal; the greedy
arrangement was proved optimal by Zalgaller and Los'~\cite{ZalgallerLos} (1994). We prove nothing about (ii). The numerical script \path{code/check_malfatti.py}
reproduces the equilateral counterexample (side~1: Malfatti circles of total area $0.315670$,
incircle plus two corner circles of radius $r/3$ of total area $0.319977$), but that is a
floating-point check, not a theorem.\footnote{The script and its output
(\path{code/check.out}) are in the malfatti repository.}

\paragraph{Statement.} For a triangle $t$ write $r$ for its inradius, $I$ for its incentre,
$s$ for its semiperimeter (\decl{Malfatti.semi}) and $a_i=|A_{i+1}A_{i+2}|$ for the side opposite
vertex $A_i$, indices in $\mathrm{Fin}\ 3$.

\begin{theorem}[\texttt{Malfatti.malfatti\_circles\_exist}, \texttt{Malfatti.malfatti\_circles\_formula}, \src{malfatti}{Malfatti}{Main}{15}, \src{malfatti}{Malfatti}{Main}{31}]\label{thm:malfatti}
Let $t$ be a triangle in $\mathrm{EuclideanSpace}\ \RR\ (\mathrm{Fin}\ 2)$. There are spheres
(circles) $c_0,c_1,c_2$ such that
\begin{itemize}
\item each $c_i$ has positive radius;
\item for pairwise distinct $i,j,k$, $c_i$ is tangent (\mdecl{IsTangent}) to the line
$A_iA_j$, and its centre lies strictly on the same side of that line as $A_k$
(\mdecl{SSameSide});
\item for $i\ne j$, $c_i$ and $c_j$ are externally tangent (\mdecl{IsExtTangent}).
\end{itemize}
Moreover the circles can be chosen with radii
\[
\rho_i=\frac{r}{2(s-a_i)}\Bigl(s-r-\bigl(|IA_{i+1}|+|IA_{i+2}|-|IA_i|\bigr)\Bigr).
\]
\end{theorem}

\begin{lstlisting}
theorem malfatti_circles_exist (t : Triangle ℝ (EuclideanSpace ℝ (Fin 2))) :
    ∃ c : Fin 3 → Sphere (EuclideanSpace ℝ (Fin 2)),
      (∀ i, 0 < (c i).radius) ∧
      (∀ i j k, i ≠ j → i ≠ k → j ≠ k →
        (c i).IsTangent line[ℝ, t.points i, t.points j] ∧
        line[ℝ, t.points i, t.points j].SSameSide (c i).center (t.points k)) ∧
      (∀ i j, i ≠ j → (c i).IsExtTangent (c j))
\end{lstlisting}
\decl{Malfatti.malfatti_circles_formula} has the same conclusion with the radius equation as an
additional first conjunct. Non-degeneracy is again free, since an \mdecl{Affine.Triangle} is
affinely independent. The general-space versions, for a triangle in any real inner product space
with \lean{Fact (finrank ℝ V = 2)}, are \decl{Malfatti.malfattiCircle_tangent_side} (which also
names the tangency point, $\mathrm{lineMap}(A_i,T,\mu)$ with $T$ the incircle touch point and
$\mu\in(0,1)$, so that it lies on the side segment), \decl{Malfatti.malfattiCircle_isExtTangent}
and \decl{Malfatti.malfattiCircle_radius_formula}; \decl{Malfatti.malfattiCircle_radius_lt_inradius}
shows $\rho_i<r$.

\paragraph{Construction and certificate.} Each circle is the image of the incircle under the
homothety at a vertex: the centre is $O_i=A_i+\mu_i(I-A_i)$ and the radius $\mu_i r$
(\decl{Malfatti.malfattiCircle}), so tangency to the two sides through $A_i$ is inherited from the
incircle (\decl{Malfatti.isTangentAt_side}). With $p,q,w$ the tangents of the quarter angles at
$A_i,A_{i+1},A_{i+2}$, defined metrically as $p=(|A_iI|-(s-a_i))/r$ (\decl{Malfatti.tparam}), the
ratio is $\mu_i=(1+q)(1+w)/(2(1+p))$ (\decl{Malfatti.mu}). The tangent length and the distance
to the incentre are $X(r,p)=r(1-p^2)/(2p)$ and $D(r,p)=r(1+p^2)/(2p)$. The one geometric input
is Heron's relation in the form $r^2s=(s-a)(s-b)(s-c)$ (\decl{Malfatti.heron}), proved from the
barycentric description of the incentre and the Gram matrix
$\langle A_i-I,A_j-I\rangle=r^2-x_ix_j$ by a \lean{linear\_combination} certificate. Heron's
relation forces the polynomial relation
\[
K(p,q,w)=1-p-q-w-pq-qw-wp+pqw=0
\]
(\decl{Malfatti.K_eq_zero_of_heron}), the tangent form of $A/4+B/4+C/4=\pi/4$; the sign of the
cofactor is controlled by $0<p,q,w<1$. Pairwise tangency (\decl{Malfatti.dist_identity}) and the
radius formula (\decl{Malfatti.radius_eq}) are then polynomial identities whose difference sides
are explicit multiples of $K$; the multipliers were found with SymPy.

\paragraph{Scope.} Only existence is proved, by verifying an explicit construction;
uniqueness of the Malfatti circles is not. Containment is proved only in part: the tangency
points lie on the side segments and the centres lie in the interior half-planes, but we do not
prove that each circle lies inside the triangle (Section~\ref{sec:limits}).

\paragraph{Numerical check.} \path{code/check_malfatti.py} computes the formula radii with centres
on the bisectors for 2{,}000 random triangles: the maximum tangency error (to the sides and
pairwise) is $2.6\times10^{-10}$, while adding $0.01$ inside the bracket of the formula gives errors
of at least $6.5\times10^{-3}$.

\section{Does the statement mean what it says?}\label{sec:correct}

\paragraph{Feuerbach.} The statement uses only Mathlib's definitions of the triangle centres, the
in-, ex- and nine-point spheres, and tangency; there is nothing local to trust. The two local
definitions of the algebra layer are bridged in \path{Feuerbach/Bridge.lean}:
\decl{Feuerbach.qf_eq_norm_sq} shows $\mathrm{qf}(\|v-u\|,\|v\|,\|u\|,p,q)=\|pu+qv\|^2$, and
\decl{Feuerbach.heron_eq_sq_two_mul_area} shows that $\mathrm{heron}$ of the three side lengths is
$(c\,b\sin\angle p_1p_0p_2)^2$. The witness triangle \decl{Feuerbach.tri345} is identified by
\decl{Feuerbach.tri345_spec} as the 3-4-5 right triangle.

\paragraph{Miquel.} The statement uses \mdecl{dist} and the local $\mathrm{Dp}$.
\path{Miquel/Bridge.lean} has 26 bridge lemmas covering all 25 local definitions: $\mathrm{Dp}$ is
the orientation determinant and vanishes exactly on \mdecl{Collinear} triples; $Q$ is the
$4\times4$ concyclicity determinant and equals $\mathrm{Dp}(P,R,S)$ times Mathlib's
\mdecl{EuclideanGeometry.Sphere.power} of the fourth point with respect to the circle through
$P,R,S$ (\decl{Miquel.Q_eq_Dp_mul_power}); the frame coefficients are identified with orientations
and powers of points; the certificate multipliers, which have no standard meaning, have bridges
recording their role; and \decl{Miquel.rpx_rpy_spec} shows that $(\mathrm{rpx},\mathrm{rpy})$ is
the only common zero of the two circle polynomials other than the origin. The docstring of
\decl{Miquel.Q} (``vanishes iff concyclic or collinear'') was misread by one automated reviewer as
a false claim; the intended reading, which is correct, is ``all four points lie on a common circle
or all four lie on a common line'', and the wording could be tightened in a later version.

\paragraph{Brahmagupta.} \decl{Brahmagupta.signedArea_spec} shows that the signed area is half
the $2\times2$ determinant of the diagonal vectors and splits along $AC$ into the signed areas of
$ABC$ and $ACD$; \decl{Brahmagupta.area_eq_half_mul_sin_angle} shows
$K=\tfrac12\,AC\cdot BD\cdot\sin\theta$ with $\theta$ the angle between the diagonals. Neither
identifies $K$ with a measure (Section~\ref{sec:limits}).

\paragraph{Malfatti.} The headline statements use only Mathlib's triangle, spheres, lines,
tangency predicates, \mdecl{SSameSide}, incentre and inradius, and the local semiperimeter.
\path{Malfatti/Bridge.lean} has nine bridge lemmas, one for each of the nine local definitions:
\decl{Malfatti.X_bridge} ($X(r,\tan\theta)=r/\tan2\theta$, the tangent length $r\cot(A/2)$);
\decl{Malfatti.D_bridge} ($D(r,\tan\theta)=r/\sin2\theta$, the distance from a vertex to the
incentre); \decl{Malfatti.K_bridge} ($K(\tan a,\tan b,\tan c)=0\iff a+b+c=\pi/4$);
\decl{Malfatti.semi_bridge} (half the perimeter, and the sum of the tangent lengths);
\decl{Malfatti.tl_bridge} ($s-a_i$, and the distance from the vertex to the incircle touch point);
\decl{Malfatti.tparam_bridge} ($\tan(A_i/4)$ with Mathlib's unoriented angle);
\decl{Malfatti.mu_bridge} ($r\mu$ at the quarter-angle tangents is Malfatti's classical radius);
\decl{Malfatti.ratio_bridge} (the ratio $\rho_i/r$, and the homothety ratio at the vertex); and
\decl{Malfatti.malfattiCircle_bridge} (the defining properties of a Malfatti circle). The bridges
hold under stated hypotheses, which we record exactly: \decl{Malfatti.K_bridge} needs
$a,b,c\in(0,\pi/4)$; \decl{Malfatti.D_bridge} needs $\theta\in(0,\pi/2)$; and the geometric bridges
(\lean{semi}, \lean{tl}, \lean{tparam}, \lean{mu}, \lean{ratio}, \lean{malfattiCircle}) assume a
2-dimensional plane, \lean{Fact (finrank ℝ V = 2)}. Only \decl{Malfatti.X_bridge} is
unconditional. These hypotheses are the natural ranges (quarter angles of a triangle lie in
$(0,\pi/4)$), but a reader should not read the bridges as unconditional identities.

\section{Related work}\label{sec:related}

Each project began with a prior-art search of Mathlib (sources and the lists \path{docs/100.yaml}
and \path{docs/1000.yaml}), GitHub code, repository and pull-request search, and other systems.
The searches are recorded in the \path{SPIKE_REPORT.md} of each repository.

\paragraph{Feuerbach: not the first Lean formalisation.} Feuerbach's theorem has earlier Lean
formalisations, and our contribution is \emph{not} priority. Weiyi Wang (GitHub \texttt{wwylele}),
the author of Mathlib's \path{NinePointCircle.lean}, has a standalone repository
\path{wwylele/Feuerbach}~\cite{WangRepo}, created on 29 January 2026, with continuous integration
passing and no \lean{sorry} in its 741-line \path{Feuerbach/Basic.lean}. It targets Lean
\texttt{v4.28.0-rc1}, follows a vector proof after Scheer~\cite{Scheer} using areas, heights,
Heron's formula and sines, and proves the internal tangency of the insphere, the external tangency
of the exspheres (for every sign set other than $\emptyset$ and the full set), the distance
$OI$ and the inequality $2r\le R$. Wang has also opened Mathlib pull request
\#44236~\cite{PR44236}, ``feat(Geometry/Euclidean): prove Feuerbach's theorem'', created on
26 September 2026; on 7 October 2026 it was open and labelled \texttt{blocked-by-other-PR} and
\texttt{merge-conflict}. At the pinned Mathlib revision (\texttt{v4.33.1}), \path{docs/100.yaml}
has no entry for theorem~29. Outside Lean, Wiedijk's list~\cite{Wiedijk100} records formalisations
in HOL~Light (Harrison), Isabelle (Paulson) and Rocq (the CertiGeo team). Our development was
written without consulting Wang's code; we found his repository and pull request only in a final
prior-art search, after our proof was complete. It is a second, independent Lean proof by a
different route (no angles, sines or Heron import; certificate-driven; about 490 lines for the
mathematical content before witnesses and bridges were added), and it should be read as such.
Given that \#44236 is open, contributing to or comparing with that pull request is the natural
route to Mathlib, rather than a competing one.

\paragraph{Independent formalisations in Atlas.} Independent Lean
formalisations of three of the four theorems were made public in the repository
\path{facebookresearch/atlas-lean}~\cite{Atlas}: \path{FeuerbachTheorem.lean},
\path{MiquelTheorem.lean} and \path{BrahmaguptaArea.lean}. All three first appeared on
5 October 2026, between 21:37 and 21:40 UTC, in a batch of import pull requests (\#46, \#61, \#62,
\#63 and \#65; the three files themselves are in \#62 and \#63). Our Miquel, Malfatti, Feuerbach
and Brahmagupta repositories were made public only on 7 October 2026, so \emph{on the public
record the Atlas files are earlier than ours}. Our own evidence of earlier completion is private:
\emph{working-copy file timestamps} (not commit timestamps) and our run log, which show the
Feuerbach algebra file on 28 September 2026 and the Miquel and Brahmagupta certificate and main
files complete by 03:45 UTC on 5 October 2026. That is, our working copies were complete before
the Atlas import, but our repositories were published after it. We did not consult the Atlas
files, and our developments are independent derivations; we claim \emph{no public priority} over
them, only independent derivation. (Atlas has no Malfatti file; our claim for Malfatti in the
Summary rests on the searches below, not on dates.) The Atlas files are imported from an upstream repository,
\path{fairinternal/formal-math}, which is private; the date on which they were first written
there is not public, so their earlier private authorship cannot be excluded either.

\paragraph{Miquel: not the first formalisation.} An earlier search by the Miquel project (Mathlib
sources and both lists, GitHub, the Archive of Formal Proofs, HOL~Light, Coq/Rocq and Mizar) found
no formal proof, but it was not repeated before writing and it missed two. The Isabelle Archive of
Formal Proofs has the entry \emph{Miquel's Theorem in Isabelle/HOL} by A.~Freitas Ramos,
D.~Barros Hulak and R.~J.~G.~B.~de~Queiroz, dated 11 August 2026~\cite{AFPMiquel}. It is
sorry-free, works in complex coordinates with the points on the side \emph{lines} (the extended
sides, as here), and proves both a pivot form (\lean{miquel\_pivot}) and a concurrence form
(\lean{miquel\_concurrent}) in which the common point is constructed, as the reflection of $r$ in
the line of centres. Mathematically this is the content of \decl{Miquel.miquel} and
\decl{Miquel.miquel_concyclic}, two months earlier. Coq's HighSchoolGeometry library (F.~Guilhot)
has \lean{Theorem Miquel}, in pivot form, proved with \lean{Qed}
(\path{theories/applications_cocyclicite.v}, line~270)~\cite{HSG}; it is relative to an
axiomatised synthetic geometry, but it is a library theorem, not merely an automated-prover
benchmark (Miquel's theorem is also a standard benchmark for Wu and Gr\"obner-basis
provers~\cite{Chou}). We therefore claim no priority for Miquel's theorem. The remaining
difference is \emph{hypothesis economy}. Our \decl{Miquel.miquel} assumes only
$\mathrm{Dp}(A,B,C)\ne0$ and $n,t,u\notin\{0,1\}$, which say exactly that the triangle is
non-degenerate and that no side point is a vertex. The AFP \lean{miquel\_concurrent} assumes these
conditions too, and in addition that the two given circles have positive radii, $0<R_1$ and
$0<R_2$, and that the two centres and the common point $r$ are not collinear,
$\lnot\,\mathrm{collinear}\ o_1\ o_2\ r$ (which excludes the case where the two circles are
tangent at $r$); its \lean{miquel\_pivot} assumes $m\ne p$, $m\ne q$, $m\ne r$, so that the
Miquel point is distinct from the three side points. Our proof handles the case $M=F$, which includes
the case where the first two circles are tangent at $F$, by the identity \decl{Miquel.tang}, and
it does not require $M$ to differ from the side points, so neither kind of extra hypothesis is
needed.

\paragraph{Brahmagupta.} In the pinned Mathlib the only \lean{brahmagupta} declaration is the
Brahmagupta--Fibonacci identity; \path{docs/1000.yaml} has an entry ``Brahmagupta theorem'' that
refers to the perpendicular-diagonals theorem, not the area formula, with no declaration; there is
no Bretschneider formula and no quadrilateral area. Mathlib does have Ptolemy's theorem and
Ptolemy's inequality, on which we build. GitHub searches (Lean, Coq, Isabelle, OCaml; \lean{brahmagupta},
\lean{Bretschneider}, ``cyclic quadrilateral'', repositories and Mathlib pull requests) found no proof
of the area formula, but that search was wrong. The ShadowBench autoformalisation
benchmark~\cite{ShadowBench} contains, under
\path{autformalization/problems/geometry/L2/geo_gen_L2_007/}, both skeleton files
\path{docs/skeletons/Skeleton1.lean}--\path{Skeleton4.lean}, which end in \lean{sorry}, and a
\emph{complete} proof in \path{ShadowBench/Source/Main.lean}~\cite{ShadowBenchProof}. At commit
\texttt{7541fcecef97cc124f38355a2cd2ba0ec8485dd5} (Lazar Milikic, 12 June 2026) that file has
334 lines, no \lean{sorry}, the main theorem \lean{brahmagupta\_formula} (line~319) and the lemma
\lean{brahmagupta\_area\_sq\_semiperimeter} (line~290). Its statement is ours: a convex cyclic
quadrilateral, the semiperimeter, $K=\sqrt{(s-a)(s-b)(s-c)(s-d)}$, with the area supplied by a
bridge hypothesis \lean{h\_area : K = quadrilateralArea \ldots}. An earlier draft of this paper
said that the closest ShadowBench results were \lean{sorry} statements; that was false, because
we had read only the skeletons. We claim no priority for Brahmagupta's formula. We have not
compiled the ShadowBench proof (Section~\ref{sec:brahmagupta}). A further search found no proof
of the all-four-point statements \decl{Brahmagupta.area_sq_le} and
\decl{Brahmagupta.area_sq_eq_iff_ptolemy}. A web and AFP search found no Brahmagupta area formula
in Isabelle, HOL~Light (whose \path{100/heron.ml} has Heron's formula only), Coq/Rocq or Mizar.
These searches are best-effort only.

\paragraph{Malfatti.} The Malfatti project searched, on 7 October 2026, the pinned Mathlib
(including \path{Archive/}, \path{docs/100.yaml} and \path{docs/1000.yaml}), the Isabelle Archive of
Formal Proofs (a mirror of 1{,}045 entries, and the AFP site search), HOL~Light, HOL4, GeoCoq, the
Rocq 100-theorems collection, the Mizar Mathematical Library, GitHub code and repository search
for Lean, Coq, Isabelle, SML and OCaml, and Wiedijk's 100 and 1000 lists, and found no
formalisation of the Malfatti circles. The only related item is an open wishlist entry in a
Mathlib issue (\#6091) for the different statement that the Malfatti circles are never optimal
(item~(ii) of Section~\ref{sec:malfatti}), not a proof. The search did not cover the Lean Zulip,
and GitHub code search indexes default branches only.

\paragraph{Summary.} The Atlas files of 5 October 2026 (above) are independent formalisations, public before our repositories,
of Feuerbach, Miquel and Brahmagupta; we claim only independent derivation, not public priority, against them. They do not change the conclusions below. As far as these searches go, we found no earlier formalisation of the
Malfatti circle construction; this is the only one of the four theorems we believe to be formalised
here for the first time. Our Brahmagupta development is a second Lean proof of a theorem already
proved in Lean in the ShadowBench repository; only its all-four-point inequality and Ptolemy
equivalence were not found elsewhere. Our Feuerbach development is
an independent formalisation of a theorem already formalised in Lean (outside Mathlib, with a
Mathlib pull request open) and in three other systems. Our Miquel development is a further
formalisation of a theorem already proved in Isabelle/HOL and in Coq, distinguished by its weaker
hypotheses. The searches are not exhaustive; both the Feuerbach and the Miquel cases show that
search can miss existing work (GitHub code search returned no hits for Wang's repository, the
first Miquel search missed an AFP entry, and the first Brahmagupta search saw only the
\lean{sorry} skeletons of a benchmark problem whose complete proof sat in the same directory).
The claim for Malfatti should be read with this record in mind.

\section{Verification}\label{sec:verify}

\emph{The Feuerbach, Miquel and Brahmagupta results in this section were obtained on 6--7~October
2026, and the Malfatti results on 7~October 2026, on the build machine (x86\_64, Linux); some were
repeated on a second machine (arm64, macOS), as stated. The working copies are not git checkouts.
For all four developments every \path{.lean} file at the published commits
(Section~\ref{sec:avail}) is byte-identical to its counterpart in the checked copies (SHA-256 of every \path{.lean} file, and \lean{cmp} on fresh clones).}

\paragraph{What the checks taught us.} We lead with the lessons that seem reusable beyond these
four libraries; the counts follow, and should be read in their light.
\begin{enumerate}
\item \emph{A witness must exercise the certificate, not merely type-check.} The first
\decl{Miquel.core} witness took the zero polynomial as its third circle; every check passed, yet it
exercised none of the certificate multipliers. It was found by reading and replaced by three
genuine circles with a second intersection point proved to be $(1,1)$, not the origin. Gate, tamper
and satisfiability layers all accepted the degenerate witness; none of them can detect this.
\item \emph{A witness should apply the theorem, so that downstream code polices it.} Rewriting
the Malfatti witnesses to state the target lemma's conclusion and prove it by applying the lemma
killed 12 of 16 genuine mutation survivors, and in Miquel the four survivors in
\path{Witness.lean} are rejected by the declarations that consume them. Witnesses written this way
are also what the harness can link to their theorems (the satisfiability table below).
\item \emph{Ablation is informative when a proof survives, not when it breaks.} The surviving
deletions found five redundant hypotheses (\lean{\_hN} in \decl{Miquel.rp_on} and four in
Malfatti), which is a concrete defect in the statements. A broken proof after deletion says only
that this proof used the hypothesis.
\end{enumerate}

\paragraph{What the battery is and is not.} The checks below are not thirteen independent
safeguards. They separate into four kinds of evidence. \emph{Kernel certification} (build,
forbidden-token grep, \lean{\#print axioms}, gate) shows that the stated propositions are proved
from the standard axioms; the tamper test checks that the gate itself rejects planted defects.
\emph{Statement-meaning evidence} (definition bridges, hand-written witnesses, numerical scripts,
and the author's reading of Table~\ref{tab:thms}) bears on whether the propositions are the
intended ones; only the bridges and witnesses are kernel-checked, and their selection is manual.
\emph{Heuristic testing} (mutation, contradiction probes, ablation, satisfiability, bridge
inventory) is incomplete by construction: a contradiction probe that finds nothing is not a
satisfiability proof, mutants are generated only for some declarations, and timeouts are counted
with type errors. \emph{Platform reproducibility} (the arm64 rebuild) checks the toolchain, not
the mathematics. The harness also has visible limits: it parsed a docstring as a declaration and
generated eight phantom mutants and a phantom satisfiability success from it, and its bridge
inventory had to be corrected by hand. Many layers share the same sources and the same author, and
their failures are correlated.

\paragraph{Build and axioms.} A clean \lean{lake build} of each library completes with 0 errors and
0 warnings; for Feuerbach, Miquel and Brahmagupta this took 8{,}710, 8{,}711 and 8{,}710 jobs
before the witness and bridge files were added (on the final sources: 8{,}712, 8{,}713, 8{,}711, and
8{,}712 for Malfatti). A
grep of the library sources for \lean{sorry}, \lean{admit}, \lean{native\_decide}, \lean{axiom},
\lean{\#eval}, \lean{run\_cmd}, \lean{set\_option}, \lean{macro}, \lean{elab}, \lean{syntax},
\lean{notation} and \lean{import Lean} is empty, and \lean{\#print axioms} reports only
\lean{propext}, \lean{Classical.choice} and \lean{Quot.sound}.

\paragraph{House battery.} The battery of the earlier AFM papers was run serially on the build
machine after the definition bridges and witnesses were added.
\begin{itemize}
\item \emph{Gate.} \path{gate.sh} rejects the forbidden tokens, builds the library, generates its
own \lean{\#print axioms} query for every declaration of the library (theorems, lemmas and
definitions), and fails on any non-standard axiom.
Result: PASS, standard axioms only, at \textbf{46} declarations (Feuerbach), \textbf{84} (Miquel,
after the \lean{core} witness was replaced, below), \textbf{30} (Brahmagupta) and \textbf{77}
(Malfatti).
\item \emph{Tamper tests.} \path{tests/tamper.sh} plants six defects in scratch copies (a
\lean{sorry}; a split \lean{axiom}; a fake \lean{\#eval} report; a macro; a deleted hypothesis; a
wrapped extra axiom) and requires each to be rejected. The deleted hypotheses are \lean{hpos} of
\decl{Feuerbach.core}, \lean{hn0} of \decl{Miquel.miquel}, \lean{hcyc} of
\decl{Brahmagupta.brahmagupta_sq} and \lean{hK} ($K=0$) of \decl{Malfatti.mu_lt_one}.
Result: \textbf{6/6} rejected in each of the four repositories.
\item \emph{Mutations.} Small edits to conclusions and hypotheses, with the proof kept, must fail
to compile. Results: Feuerbach 177/177 rejected, no survivors; Brahmagupta 78/84 rejected,
6 survivors; Miquel 192/199 rejected, 7 survivors; Malfatti 226/238 rejected, 12 survivors, of
which 4 are genuine (below). ``Rejected'' includes failures by heartbeat timeout (3 for Feuerbach,
8 for Miquel, 1 for Malfatti). These totals cover only the declarations for which the harness
generated mutants, and several headline statements are not among them. For Feuerbach the harness
generated \emph{no} mutants for \decl{Feuerbach.isIntTangent_insphere_ninePointCircle},
\decl{Feuerbach.isExtTangent_exsphere_ninePointCircle} and \decl{Feuerbach.nontrivial_of_triangle},
so ``177/177 rejected'' says nothing about the two tangency theorems. For Brahmagupta it generated
none for \decl{Brahmagupta.brahmagupta}, \decl{Brahmagupta.brahmagupta_sq},
\decl{Brahmagupta.area_sq_le}, \decl{Brahmagupta.area_sq_eq_iff_ptolemy},
\decl{Brahmagupta.brahmagupta_identity} and \decl{Brahmagupta.not_cospherical_of_area_sq_ne},
that is, for every headline statement of Section~\ref{sec:brahmagupta}; the 84 Brahmagupta
mutants all lie in the auxiliary lemmas and examples. The five Malfatti declarations without
mutants are listed below.
All six Brahmagupta survivors lie inside the concrete examples and were reproduced and classified
harmless by hand (for instance, a ``not concyclic'' assertion survives almost any small nudge of a
point, because landing exactly on a circle is non-generic).
\emph{Miquel.} All Miquel battery figures in this section are from the run on the current
sources completed on 7 October 2026 (started 17:58 UTC), after the \lean{core} witness was
replaced. There are seven survivors. Three are in
\decl{Miquel.miquel_negative_control} (\path{Miquel/Main.lean}) and four in
\path{Miquel/Witness.lean} (two in \decl{Miquel.miquel_hyps_example}, two in
\decl{Miquel.core_example}). All seven were reproduced (each mutant compiles with the original
proof) and ruled on by hand: all seven are harmless and none is a genuine defect. Each mutated
statement is still a true closed numerical fact under an open ($\ne0$) condition; in the two
\decl{Miquel.core_example} slots the side condition reduces to $4+4a^2$ and $b^2+4$, strictly
positive for every real value. For the four survivors in \path{Witness.lean} the downstream
declaration that consumes the witness (\decl{Miquel.miquel_point_example}, or
\decl{Miquel.core_example_conclusion} through \decl{Miquel.core}) rejects the mutant with a type
mismatch.
\emph{Malfatti.} The non-vacuity witnesses were rewritten so that each states the target lemma's
conclusion at concrete arguments and is proved by applying the lemma, with its hypotheses
discharged inline; this killed 12 of the 16 genuine survivors of the previous run. Of the 12
survivors in the final run, 8 are an artefact of the harness, which parsed a docstring as a
declaration. The 4 that remain are benign: the mutants of the library lemmas
\decl{Malfatti.fin3_distinct} (one) and \decl{Malfatti.fin3_third} (two) are true but weaker
statements about $\mathrm{Fin}\ 3$, killed downstream where the dropped conjunct is used, and the
mutant of \decl{Malfatti.heron_witness} ($1^2\mapsto1^3$) is value-preserving, since $1^3=1^2$.
For five declarations (\decl{Malfatti.isTangentAt_side}, \decl{Malfatti.malfattiCircle_isExtTangent},
\decl{Malfatti.malfattiCircle_tangent_side}, \decl{Malfatti.sSameSide_center},
\decl{Malfatti.touch_mem}) the harness generated no mutants at all (Section~\ref{sec:limits}).
The general lesson is that mutation testing is weak on concrete numerical examples and on
non-generic negative claims.
\item \emph{Contradictory hypotheses.} No hypothesis set was found contradictory (0 in each).
\item \emph{Ablation.} Each hypothesis is deleted in turn and the unchanged proof is recompiled.
Result, as (proof broke)/(proof survived): Feuerbach 70/0; Miquel 186/1; Brahmagupta 6/0; Malfatti
66/4. The harness labels the first figure ``needed'', but that is too strong: a broken proof shows
only that \emph{this} proof uses the hypothesis, not that the statement is false without it. The
second figure is the informative one, since a surviving proof does show the hypothesis redundant. The unused Miquel hypothesis is \lean{\_hN} of \decl{Miquel.rp_on}, already
named with an underscore. It is genuinely unnecessary: if $\mathrm{rpN}=0$ then, with Lean's
convention $x/0=0$, the point $(\mathrm{rpx},\mathrm{rpy})$ is the origin, where both circle
polynomials vanish. The four unused Malfatti hypotheses are $0<r$ in \decl{Malfatti.radius_eq},
$p<1$ in \decl{Malfatti.dist_identity}, $q<1$ in \decl{Malfatti.mu_lt_one} and $i\ne j$ in
\decl{Malfatti.fin3_third} (Section~\ref{sec:limits}).
\item \emph{Unused-variable linter.} 0 warnings in all four libraries.
\item \emph{Satisfiability.} For each theorem declaration, the harness tries to produce a
kernel-checked \lean{example} exhibiting values for all its hypotheses. Since the reviewers'
reports, the harness has a separate route for hand-written witnesses: a witness declaration that
applies a theorem with its hypotheses discharged is \emph{linked} to that theorem and reported on
its own line, never folded into the certified count. The denominator is \emph{all} theorem
declarations, including hypothesis-free and closed statements, which are certified trivially; the
ratios below are therefore not coverage of hypothesis-bearing statements. Results, separated:
\begin{center}\small
\begin{tabular}{@{}lrrrr@{}}
\toprule
 & machine-certified & witness-linked & unresolved & total \\ \midrule
Feuerbach & 30 & 0 & \textbf{13} & 43 \\
Miquel & 40 & 0 & \textbf{19} & 59 \\
Brahmagupta & 22 & 0 & \textbf{4} & 26 \\
Malfatti & 26 & 33 & \textbf{8} & 67 \\
\bottomrule
\end{tabular}
\end{center}
For Malfatti, $26+33=59$ of 67 declarations are credited and 8 are unresolved; the total is 67, not
the 68 of an earlier run, because the harness no longer parses a docstring as a declaration. For the
other three libraries no witness is written in the form the linker recognises, so the
witness-linked count is~0 and the unresolved counts stand as reported: \textbf{13} Feuerbach items
(algebra and frame lemmas), \textbf{19} Miquel items (bridge and specification lemmas together with
\decl{Miquel.dist_eq_of_Q_eq_zero}, \decl{Miquel.circ_aux}, \decl{Miquel.core},
\decl{Miquel.bridge}, \decl{Miquel.miquel} and \decl{Miquel.miquel_concyclic}) and \textbf{4}
Brahmagupta items. Several of these, including every headline statement, are exercised by the
hand-written witnesses listed below, but that coverage is established by reading, not by the
harness, and we do not count it. \decl{Miquel.rp_zero} is machine-certified. An unresolved item is
one for which no satisfying instance was produced automatically; it is not evidence of vacuity, and
a certified item is evidence only that the hypotheses are jointly satisfiable.
\item \emph{Bridge inventory.} The scanner's heuristic reports Miquel 21 of 25 notions bridged,
Brahmagupta 2 of 4, Feuerbach 0 of 3 and Malfatti 1 of 10. These numbers are under-counts: the
heuristic does not match the names in the \path{Bridge.lean} files. The inventory we rely on was
therefore corrected by hand: by inspection every mathematical local definition of the four
libraries has a bridge lemma (Section~\ref{sec:correct}); for Malfatti that is nine bridges for
nine definitions, the tenth definition being the witness triangle. This is a reading, not a
machine check.
\end{itemize}

\paragraph{Non-vacuity witnesses.} \emph{Feuerbach:} \decl{Feuerbach.feuerbach_nonvacuous}
instantiates every hypothesis of \decl{Feuerbach.feuerbach} with $V=P=\CC$ (of real dimension~2)
and the triangle \decl{Feuerbach.tri345} with vertices $0,4,3i$, whose affine independence is
\decl{Feuerbach.tri345_affineIndependent}; \decl{Feuerbach.algebra_hyps_345} shows that the
hypotheses of the algebra lemmas are jointly satisfiable ($a,b,c=5,3,4$, $R=5/2$, $k=12$), and
\decl{Feuerbach.nine_in_345} instantiates \decl{Feuerbach.nine_in} there.
\emph{Miquel:} \decl{Miquel.miquel_hyps_example} discharges all ten hypotheses of
\decl{Miquel.miquel} for the midpoint configuration, and \decl{Miquel.miquel_point_example}
applies the theorem there; further witnesses cover \decl{Miquel.dist_eq_of_Q_eq_zero},
\decl{Miquel.circ_aux}, \decl{Miquel.bridge} and \decl{Miquel.core}. The \decl{Miquel.core}
witness \decl{Miquel.core_example} uses three genuine circles, $x^2+y^2-2x$, $x^2+y^2-2y$ (meeting
at the origin and at $(1,1)$) and $x^2+y^2-2$, with non-zero certificate multipliers;
\decl{Miquel.core_example_point} proves that the second intersection point of the first two is
$(1,1)$, not the origin, and \decl{Miquel.core_example_conclusion} applies \decl{Miquel.core} to put
it on the third circle. (An earlier version of this witness took the zero polynomial as its third
circle and so did not exercise the certificates; it has been replaced.) \decl{Miquel.rp_zero} has
no hand-written witness; it is harness-certified satisfiable.
\emph{Brahmagupta:} the unit square and the rhombus of Section~\ref{sec:brahmagupta}, through
\decl{Brahmagupta.unitSquare_brahmagupta}, \decl{Brahmagupta.unitSquare_square_area} and
\decl{Brahmagupta.rhombus_not_cospherical'}.
\emph{Malfatti:} \decl{Malfatti.malfatti_circles_exist_triW} and
\decl{Malfatti.malfatti_circles_formula_triW} instantiate the two headline theorems at the right
triangle \decl{Malfatti.triW} with vertices $(0,0),(1,0),(0,1)$, whose affine independence is
\decl{Malfatti.triW_affineIndependent}. The algebraic lemmas are instantiated at $r=1$,
$(p,q,w)=(1/3,1/3,1/7)$ and, so that the $r$-scaling is exercised, at $r=2$,
$(p,q,w)=(1/3,1/5,3/11)$: Heron's relation (\decl{Malfatti.heron_witness},
\decl{Malfatti.heron_witness_r2}) is fed to \decl{Malfatti.K_eq_zero_of_heron}, whose conclusion
$K=0$ is in turn fed to \decl{Malfatti.radius_eq}, \decl{Malfatti.dist_identity} and
\decl{Malfatti.mu_lt_one}; \decl{Malfatti.param_spec} is instantiated at $(r,x,d)=(3,4,5)$. Each
witness is the target lemma's conclusion at those arguments, so most changes to the lemma or to a
numeral break the application (not all: see the value-preserving survivor above).

\paragraph{Cross-platform rebuild on arm64/macOS.} All four libraries
were rebuilt on a second machine with a different architecture and operating system (Apple
M5~Pro, arm64, macOS), with the same Lean and Mathlib versions, from sources transferred by
tarball and checked by SHA-256 on both legs. Before the witnesses and bridges were added, the
three builds there completed with exit status~0, 0 warnings and the same job counts as on the
build machine (8{,}710, 8{,}711, 8{,}710), and \lean{\#print axioms} agreed for all 25
declarations then cited. On the final sources all four libraries built there with exit status~0
and 0 warnings (8{,}712, 8{,}713, 8{,}711 and 8{,}712 jobs), \path{gate.sh} passed at 46, 84, 30
and 77 declarations, and \path{tests/tamper.sh} rejected all six defects in each repository. This
rebuild uses the same sources, toolchain and author, so it checks platform independence, not
independent authorship.

\paragraph{Limits of these checks.} None of these checks shows that a formal statement says what
its informal counterpart says. The author has checked each entry of Table~\ref{tab:thms} by hand
against the cited Lean declaration and confirms that the formal statement, with the listed
hypotheses, is the informal statement in the same row. This is a single reader's judgement, so it
was complemented by a blind re-formalisation.

\paragraph{Blind re-formalisation.} The paper's text, with every Lean name, declaration
reference and repository link removed, was given to a separate language-model writer that had no
access to the code. From that text alone it wrote a Lean statement of each of the four theorems.
Each blind statement was then compared with the corresponding library theorem by proving, in Lean
against the pinned libraries, both implications: blind statement $\Rightarrow$ library statement
and library statement $\Rightarrow$ blind statement. Inside these proofs only the premise, the
local definitions and general Mathlib lemmas were allowed; a transitive walk over the dependency
closure of all eight implications confirms that none of the fourteen headline theorems of the
libraries is used. All eight implications were proved, from the standard axioms only, so all four
pairs are equivalent. Two points are worth recording. For Feuerbach the blind statement has no
hypothesis that the ambient space is two-dimensional; the library theorem does, and the implication
from the library statement to the blind one goes through by restricting to the affine span of the
triangle, confirming that this hypothesis is removable, as stated in Section~\ref{sec:feuerbach}.
For Feuerbach and Malfatti, identifying the blind statement's explicitly defined centres and radii
with Mathlib's uses the library's coordinate lemmas for those points, which are definition bridges
and say nothing about tangency. The comparison is statement-level: it shows that the prose and the
code state the same propositions, not that either is true. The writer was a model, not an
independent human formaliser, and it shares the paper's wording; the check guards against
prose-versus-code drift, not against a shared misreading of the mathematics.

\section{Limitations and further work}\label{sec:limits}

\begin{itemize}
\item \emph{Feuerbach: prior art.} The theorem is already formalised in Lean by Wang
(standalone repository and open Mathlib pull request \#44236) and in HOL~Light, Isabelle and Rocq.
Our development is independent but not first. Smaller points: only the three excircles with sign
set $\{i\}$ are stated; the tangency point (the Feuerbach point) is not given explicitly, although
Mathlib's characterisation produces it inside an existential; and the long machine-generated
\lean{linear\_combination} certificates would need shortening before they could be offered to
Mathlib.
\item \emph{Feuerbach: a global instance.} \path{Feuerbach/Witness.lean} declares a
\lean{noncomputable instance : Inhabited (Triangle ℝ ℂ)} so that the automated satisfiability
checker has a default triangle. Adding a default instance can in principle affect elaboration
elsewhere; the paper's non-vacuity claim rests on \decl{Feuerbach.feuerbach_nonvacuous}, not on the
instance. We have decided to keep it in the published version: it adds no axiom and no hypothesis
to any theorem, the satisfiability layer uses it as its default triangle, and removing it would
change the pinned commit on which every check of Section~\ref{sec:verify} was run. It is an orphan
instance on a Mathlib type, so it would be dropped from any Mathlib submission.
\item \emph{Brahmagupta: prior art.} Brahmagupta's formula is already proved in Lean~4, without
\lean{sorry}, in the ShadowBench repository (commit of 12 June 2026). Our development is not first;
what it adds is the all-four-point inequality and Ptolemy equivalence, the statement in Mathlib's
vocabulary and the verification record. We did not compile the earlier proof.
\item \emph{Brahmagupta: area is the shoelace value.} \decl{Brahmagupta.area} is the absolute
value of the shoelace expression. It is \emph{not} identified with the Lebesgue measure of the
convex hull of $\{A,B,C,D\}$; that bridge needs polygon-measure groundwork that Mathlib does not
have. For a convex quadrilateral, which the headline theorems assume, the shoelace value is the
geometric area, but this is not proved inside Lean.
\item \emph{Brahmagupta: no trigonometric Bretschneider.} The angle form
$K^2=(s-a)(s-b)(s-c)(s-d)-abcd\cos^2\tfrac{\alpha+\gamma}2$ is not proved; only the equivalent,
angle-free diagonal form \decl{Brahmagupta.sixteen_area_sq} is.
\item \emph{Brahmagupta: degenerate configurations.} The hypotheses of
\decl{Brahmagupta.brahmagupta} (\mdecl{Cospherical} and two \mdecl{Sbtw} conditions) do not force the
four vertices to be distinct: for instance $A=B$ and $C=D$ on a diameter are allowed. In such
collapsed cases the formula holds with both sides~$0$. The informal reading ``convex cyclic
quadrilateral'' is accurate only when the vertices are also distinct (Section~\ref{sec:brahmagupta}).
\item \emph{Brahmagupta: non-convex orderings.} Crossed (self-intersecting) orderings of four
concyclic points are not covered. For them \decl{Brahmagupta.brahmagupta_sq} does not apply, and
\decl{Brahmagupta.area_sq_eq_iff_ptolemy} shows that the formula fails unless Ptolemy's equality
holds for that ordering.
\item \emph{Miquel: model and degeneracies.} The theorem is stated in $\CC$ with explicit
coordinates and \mdecl{dist}, not with Mathlib's \mdecl{EuclideanGeometry.Sphere} or
\mdecl{EuclideanGeometry.Cospherical} API; a restatement would be routine but has not been done.
Degenerate configurations (a collinear triangle, or $D,E,F$ at a vertex) are excluded rather than
handled. $M$ is not proved distinct from $D,E,F$, and can coincide with one of them in tangent
configurations. The Miquel-point corollaries (spiral similarity, the complete-quadrilateral and
pentagon theorems) are not formalised.
\item \emph{Miquel: prior art.} Miquel's theorem is already formalised in Isabelle/HOL (AFP,
August 2026) and in Coq's HighSchoolGeometry. Our development is not first; its distinguishing
feature is that it needs no hypotheses beyond non-degeneracy of the triangle and of the side points.
\item \emph{Miquel: an unwitnessed lemma and a decoupled control.} \decl{Miquel.rp_zero} has no
hand-written witness; it is harness-certified satisfiable. \decl{Miquel.miquel_negative_control} is a standalone numerical
non-concyclicity fact with no formal link to \decl{Miquel.miquel_example}; it supports the claim
that membership of the third circle is a real constraint only by inspection.
\item \emph{Malfatti: not the conjecture.} We formalise the construction and the radius formula
only. Malfatti's maximal-area conjecture is false (Goldberg 1967) and is not formalised, nor is the
optimality of the greedy arrangement (Zalgaller and Los' 1994).
\item \emph{Malfatti: uniqueness and containment.} Uniqueness of the Malfatti circles is not
proved. We do not prove that each circle lies inside the triangle, only that its tangency points
lie on the side segments and its centre in the interior half-planes.
\item \emph{Malfatti: conditional bridges.} \decl{Malfatti.K_bridge} holds only for
$a,b,c\in(0,\pi/4)$, \decl{Malfatti.D_bridge} only for $\theta\in(0,\pi/2)$, and the geometric
bridges assume a 2-dimensional plane (Section~\ref{sec:correct}).
\item \emph{Malfatti: four unused hypotheses.} Ablation finds four hypotheses that the proofs do not
use: $0<r$ in \decl{Malfatti.radius_eq}, $p<1$ in \decl{Malfatti.dist_identity}, $q<1$ in
\decl{Malfatti.mu_lt_one} and $i\ne j$ in \decl{Malfatti.fin3_third} (for any $i,j\in\mathrm{Fin}\ 3$
some $k$ avoids both). None affects soundness, but each lemma is stated more weakly than it could
be, and the hypotheses should be dropped or justified before publication.
\item \emph{Malfatti: declarations without mutants.} The mutation harness generated no mutants
for \decl{Malfatti.isTangentAt_side}, \decl{Malfatti.malfattiCircle_isExtTangent},
\decl{Malfatti.malfattiCircle_tangent_side}, \decl{Malfatti.sSameSide_center} and
\decl{Malfatti.touch_mem}, so the mutation layer says nothing about these five statements; they
have been checked only by reading and by their use in the headline theorems.
\item \emph{Feuerbach and Brahmagupta: declarations without mutants.} The harness generated no
mutants for the Feuerbach tangency theorems \decl{Feuerbach.isIntTangent_insphere_ninePointCircle}
and \decl{Feuerbach.isExtTangent_exsphere_ninePointCircle} or for
\decl{Feuerbach.nontrivial_of_triangle}, nor for Brahmagupta's headline statements
\decl{Brahmagupta.brahmagupta}, \decl{Brahmagupta.brahmagupta_sq}, \decl{Brahmagupta.area_sq_le},
\decl{Brahmagupta.area_sq_eq_iff_ptolemy}, \decl{Brahmagupta.brahmagupta_identity} and
\decl{Brahmagupta.not_cospherical_of_area_sq_ne}. The mutation layer says nothing about these nine
statements; they have been checked by reading, by the witnesses and by the numerical scripts.
\end{itemize}

\section*{Acknowledgements and use of AI}
A large-language-model AI assistant was used to write the Lean proofs and the Python checking
scripts, to run the prior-art searches and the verification battery, and to draft this paper. The
author directed the work and reviewed the paper. Every proof is checked by the Lean kernel, and
the author takes responsibility for all content.

\begin{thebibliography}{99}\sloppy
\bibitem{AFP} \emph{Archive of Formal Proofs}, \url{https://www.isa-afp.org/}.
\bibitem{AFPMiquel} A.~Freitas Ramos, D.~Barros Hulak and R.~J.~G.~B.~de~Queiroz, Miquel's
Theorem in Isabelle/HOL, \emph{Archive of Formal Proofs}, 11 August 2026,
\url{https://isa-afp.org/entries/Miquel.html}; source \path{thys/Miquel/Miquel.thy} in
\url{https://github.com/isabelle-prover/mirror-afp-devel}, accessed 7 October 2026.
\bibitem{HSG} F.~Guilhot et al., \emph{HighSchoolGeometry} (Rocq community library),
\path{theories/applications_cocyclicite.v}, \lean{Theorem Miquel},
\url{https://github.com/rocq-community/HighSchoolGeometry}, accessed 7 October 2026.
\bibitem{Atlas} \emph{atlas-lean}, \url{https://github.com/facebookresearch/atlas-lean}, files
\path{FeuerbachTheorem.lean}, \path{MiquelTheorem.lean} and \path{BrahmaguptaArea.lean}, imported
in pull requests \#46, \#61, \#62, \#63 and \#65 on 5 October 2026 (21:37--21:40 UTC) from the
private repository \path{fairinternal/formal-math}; accessed 7 October 2026.
\bibitem{Chou} S.-C.~Chou, \emph{Mechanical Geometry Theorem Proving}, D.~Reidel, Dordrecht, 1988.
\bibitem{Goldberg} M.~Goldberg, On the original Malfatti problem, \emph{Mathematics Magazine}
40 (1967), 241--247.
\bibitem{Lean4} L.~de Moura and S.~Ullrich, The Lean~4 theorem prover and programming language, in
\emph{Automated Deduction -- CADE~28}, Lecture Notes in Computer Science 12699, Springer, 2021,
625--635.
\bibitem{LobRichmond} H.~Lob and H.~W.~Richmond, On the solutions of Malfatti's problem for a
triangle, \emph{Proceedings of the London Mathematical Society} (2) 30 (1930), 287--304.
\bibitem{Mathlib} The mathlib Community, The Lean mathematical library, in \emph{Proceedings of the
9th ACM SIGPLAN International Conference on Certified Programs and Proofs (CPP 2020)}, ACM, 2020,
367--381. Repository: \url{https://github.com/leanprover-community/mathlib4}.
\bibitem{PaperA} J.~Bald, Classical plane geometry by coordinate certificates: Desargues, Monge,
Simson, Napoleon and Lami in Lean~4, manuscript, 2026, not yet published. Lean sources:
\url{https://github.com/jpbald93/desargues}, \url{https://github.com/jpbald93/monge},
\url{https://github.com/jpbald93/simson}, \url{https://github.com/jpbald93/napoleon} and
\url{https://github.com/jpbald93/lami}.
\bibitem{PR44236} W.~Wang, Mathlib pull request \#44236, \emph{feat(Geometry/Euclidean): prove
Feuerbach's theorem}, opened 26 September 2026, open at the time of access,
\url{https://github.com/leanprover-community/mathlib4/pull/44236}, accessed 7 October 2026.
\bibitem{Scheer} M.~Scheer, A simple vector proof of Feuerbach's theorem, arXiv:1107.1152.
\bibitem{ShadowBench} ShadowBench problem \path{geo_gen_L2_007},
\url{https://github.com/epfl-lara/icml-26-lean-challenges/tree/main/autformalization/problems/geometry/L2/geo_gen_L2_007};
accompanying paper \url{https://arxiv.org/html/2608.29270v1}, \S J.5.
\bibitem{ShadowBenchProof} L.~Milikic (EPFL LARA), \lean{brahmagupta\_formula}, file
\path{autformalization/problems/geometry/L2/geo_gen_L2_007/ShadowBench/Source/Main.lean} in
\url{https://github.com/epfl-lara/icml-26-lean-challenges}, commit
\texttt{7541fcecef97cc124f38355a2cd2ba0ec8485dd5}, 12 June 2026 (commit message ``Checkpoint
autoformalization proof state''), accessed 7 October 2026.
\bibitem{WangRepo} W.~Wang, \emph{Feuerbach} (Formalization of Feuerbach's Theorem),
\url{https://github.com/wwylele/Feuerbach}, created 29 January 2026, accessed 7 October 2026.
\bibitem{ZalgallerLos} V.~A.~Zalgaller and G.~A.~Los', The solution of Malfatti's problem,
\emph{Journal of Mathematical Sciences} 72 (1994), 3163--3177.
\bibitem{Wiedijk100} F.~Wiedijk, \emph{Formalizing 100 Theorems},
\url{https://www.cs.ru.nl/~freek/100/}, accessed 7 October 2026.
\end{thebibliography}

\end{document}
